\documentclass[../../main.tex]{subfiles}

\graphicspath{{../../image/}{../../figure/}}

\begin{document}

\section{习题}\label{section:ch07-习题}

\begin{exercise}\label{exercise:ch07-01}
（★）通过将\eqref{align:ch07-特征展开}代入\eqref{align:ch07-极小点展开}，并使用\eqref{align:ch07-特征方程}和\eqref{align:ch07-正交规范}，证明误差函数\eqref{align:ch07-极小点展开}可以写成\eqref{align:ch07-对角化}的形式。
\end{exercise}

\begin{exercise}\label{exercise:ch07-02}
（★）考虑具有特征向量方程\eqref{align:ch07-特征方程}的黑塞矩阵 $\mathbf{H}$。通过将\eqref{align:ch07-正定二次型展开}中的向量 $\mathbf{v}$ 依次设置为每个特征向量 $\mathbf{u}_i$，证明黑塞矩阵 $\mathbf{H}$ 是正定的——当且仅当它的所有特征值都是正的。
\end{exercise}

\begin{exercise}\label{exercise:ch07-03}
（★★）通过考虑关于驻点 $\mathbf{w}^{\ast}$ 的误差函数的局部 Taylor 展开\eqref{align:ch07-极小点展开}，证明该驻点是误差函数的局部极小值的充分必要条件是，由\eqref{align:ch07-黑塞矩阵}定义的黑塞矩阵 $\mathbf{H}$ 在 $\hat{\mathbf{w}}=\mathbf{w}^{\ast}$ 的情况下是正定的。
\end{exercise}

\begin{exercise}\label{exercise:ch07-04}
（★★）考虑一个具有单一输入变量 $x$ 和单一输出变量 $y$ 的线性回归模型：
\begin{align}
y(x,w,b) = wx+b,
\label{align:ch07-习题线性模型}
\end{align}
平方和误差函数为
\begin{align}
E(w,b) = \frac{1}{2}\sum_{n=1}^{N}\bigl\{y(x_n,w,b)-t_n\bigr\}^{2}.
\label{align:ch07-习题平方和}
\end{align}
推导出由误差函数对权重参数 $w$ 和偏置参数 $b$ 的二阶导数给出的 $2\times2$ 黑塞矩阵的元素表达式。证明这个黑塞矩阵的迹和行列式都是正的。由于矩阵的迹代表特征值之和，而行列式对应特征值的乘积（参见\cref{section:appA-迹和行列式}），两个特征值都是正的，因此误差函数在驻点达到最小值。
\end{exercise}

\begin{exercise}\label{exercise:ch07-05}
（★★）考虑一个带有单一输入变量 $x$ 和单一输出变量 $y$ 的单层分类模型：
\begin{align}
y(x,w,b) = \sigma(wx+b),
\label{align:ch07-习题分类模型}
\end{align}
其中 $\sigma(\cdot)$ 是\cref{definition:ch05-logistic-sigmoid}中定义的逻辑斯谛 sigmoid 函数。交叉熵误差函数为
\begin{align}
E(w,b) = -\sum_{n=1}^{N}\bigl\{t_n\ln y(x_n,w,b)+(1-t_n)\ln\bigl(1-y(x_n,w,b)\bigr)\bigr\}.
\label{align:ch07-习题交叉熵}
\end{align}
推导出由误差函数对权重参数 $w$ 和偏置参数 $b$ 的二阶导数给出的 $2\times2$ 黑塞矩阵的元素表达式。证明这个黑塞矩阵的迹和行列式都是正的。由于矩阵的迹代表特征值之和，而行列式对应特征值的乘积（参见\cref{section:appA-迹和行列式}），两个特征值都是正的，因此误差函数在驻点达到最小值。
\end{exercise}

\begin{exercise}\label{exercise:ch07-06}
（★★）考虑一个由\eqref{align:ch07-极小点展开}定义的二次误差函数，在该误差函数中，黑塞矩阵 $\mathbf{H}$ 有一个由\eqref{align:ch07-特征方程}给出的特征方程。证明常数误差的等高线呈椭圆形，其轴与特征向量 $\mathbf{u}_i$ 对齐，长度则与相应特征值 $\lambda_i$ 的平方根成反比。
\end{exercise}

\begin{exercise}\label{exercise:ch07-07}
（★）证明由于黑塞矩阵 $\mathbf{H}$ 的对称性，二次误差函数\eqref{align:ch07-Taylor展开}中独立元素的数量由 $W(W+3)/2$ 给出。
\end{exercise}

\begin{exercise}\label{exercise:ch07-08}
（★）考虑一组值 $x_1,\cdots,x_N$，它们是从均值为 $\mu$、方差为 $\sigma^{2}$ 的分布中抽取的，并定义样本均值为
\begin{align}
\bar{x} = \frac{1}{N}\sum_{n=1}^{N}x_n.
\label{align:ch07-习题样本均值}
\end{align}
证明关于数据来源分布的平方误差 $(\bar{x}-\mu)^{2}$ 的期望由 $\sigma^{2}/N$ 给出。这说明样本的均方根误差由 $\sigma/\sqrt{N}$ 给出，并且会随着样本数量的增加而相对缓慢地减小。
\end{exercise}

\begin{exercise}\label{exercise:ch07-09}
（★★）考虑一个在第 $l$ 层计算函数\eqref{align:ch07-前向计算}和\eqref{align:ch07-激活函数}的分层网络。假设我们使用 Gauss 分布 $\mathcal{N}(0,\varepsilon^{2})$ 初始化权重，并假设第 $(l-1)$ 层单元的输出 $z_j^{(l-1)}$ 具有方差 $\lambda^{2}$。通过使用 ReLU 激活函数的形式，证明第 $l$ 层输出的平均值和方差分别由\eqref{align:ch07-初始化均值}和\eqref{align:ch07-初始化方差}给出，进而证明如果我们想要第 $l$ 层的单元的预激活也有方差 $\lambda^{2}$，则 $\varepsilon$ 的值应该由\eqref{align:ch07-初始化尺度}给出。
\end{exercise}

\begin{exercise}\label{exercise:ch07-10}
（★★）通过使用\eqref{align:ch07-极小点展开}、\eqref{align:ch07-特征方程}和\eqref{align:ch07-特征展开}，推导出结果\eqref{align:ch07-梯度展开}和\eqref{align:ch07-权重变化}，这两个结果以黑塞矩阵的特征向量展开形式来表达梯度向量和一般的权重更新。利用这两个结果，并结合特征向量的正交性关系\eqref{align:ch07-正交规范}和梯度下降公式\eqref{align:ch07-批量梯度下降}，导出以系数 $\alpha_i$ 表示的批量梯度下降更新结果\eqref{align:ch07-alpha变化}。
\end{exercise}

\begin{exercise}\label{exercise:ch07-11}
（★）考虑一个平滑变化、曲率较低的误差表面，梯度随位置缓慢变化。证明对于较小的学习率和动量参数，由\eqref{align:ch07-Nesterov}定义的 Nesterov 动量梯度更新等价于由\eqref{align:ch07-动量更新}定义的标准带动量的梯度下降。
\end{exercise}

\begin{exercise}\label{exercise:ch07-12}
（★★）考虑变量 $x$ 的一系列值 $\{x_1,\cdots,x_N\}$，并假设我们使用下式计算指数加权移动平均值：
\begin{align}
\mu_n = \beta\mu_{n-1}+(1-\beta)x_n,
\label{align:ch07-习题ewma}
\end{align}
其中 $0\leqslant\beta\leqslant1$。通过使用有限几何级数求和的以下结果：
\begin{align}
\sum_{k=1}^{n}\beta^{k-1} = \frac{1-\beta^{n}}{1-\beta},
\label{align:ch07-习题几何级数}
\end{align}
证明如果使用 $\mu_0=0$ 初始化平均序列，则估计器存在偏差，并且可以使用下式校正偏差：
\begin{align}
\hat{\mu}_n = \frac{\mu_n}{1-\beta^{n}}.
\label{align:ch07-习题偏差校正}
\end{align}
\end{exercise}

\begin{exercise}\label{exercise:ch07-13}
（★）在梯度下降法中，权重向量 $\mathbf{w}$ 是通过在权重空间中沿由学习率参数 $\eta$ 控制的负梯度方向进行更新的。假设我们选择权重空间中的一个方向 $\mathbf{d}$，沿着该方向在给定当前权重向量 $\mathbf{w}^{(\tau)}$ 的情况下最小化误差函数
\begin{align}
E\bigl(\mathbf{w}^{(\tau)}+\lambda\mathbf{d}\bigr),
\label{align:ch07-习题线搜索}
\end{align}
进行最小化，以获得新权重向量 $\mathbf{w}^{(\tau+1)}$ 所对应的值 $\lambda^{\ast}$。证明在 $\mathbf{w}^{(\tau+1)}$ 处，$E(\mathbf{w})$ 的梯度垂直于向量 $\mathbf{d}$。这称为“线搜索”方法，它是各种数值优化算法的基础（Bishop, 1995b）。
\end{exercise}

\begin{exercise}\label{exercise:ch07-14}
（★）证明由\eqref{align:ch07-输入标准化}定义的重新归一化的输入变量，其中 $\mu_i$ 由\eqref{align:ch07-输入均值}定义，$\sigma_i^{2}$ 由\eqref{align:ch07-输入方差}定义，具有零均值和单位方差。
\end{exercise}

\end{document}
